% Simple LaTeX document - compatible with pdflatex
\documentclass[12pt,a4paper]{report}
\usepackage[margin=1in]{geometry}
\usepackage[utf8]{inputenc}
\usepackage{graphicx}
\usepackage{hyperref}
\usepackage{amsmath}
\usepackage{amssymb}
\usepackage{listings}
\usepackage{xcolor}
\usepackage{longtable}
\usepackage{booktabs}
\usepackage{caption}

% Color definitions
\definecolor{codegreen}{rgb}{0,0.6,0}
\definecolor{codegray}{rgb}{0.5,0.5,0.5}
\definecolor{codepurple}{rgb}{0.58,0,0.82}
\definecolor{backcolour}{rgb}{0.95,0.95,0.92}
\definecolor{navy}{RGB}{0,0,128}

\hypersetup{
    colorlinks=true,
    linkcolor=navy,
    filecolor=magenta,
    urlcolor=navy,
}

% Code listing style
\lstdefinestyle{mystyle}{
    backgroundcolor=\color{backcolour},
    commentstyle=\color{codegreen},
    keywordstyle=\color{blue},
    numberstyle=\tiny\color{codegray},
    stringstyle=\color{codepurple},
    basicstyle=\ttfamily\small,
    breaklines=true,
    frame=single,
    numbers=left
}
\lstset{style=mystyle}

\title{TSL-51: Thai Sign Language Recognition System}
\author{TSL Research Team}
\date{April 2026}

\begin{document}

\maketitle

\chapter{บทนำ}

\section{ความเป็นมา}

ภาษามือ (Sign Language) เป็นรูปแบบการสื่อสารหลักสำหรับผู้บกพร่องทางการได้ยิน ในประเทศไทยมีผู้ใช้ภาษามือไทย (Thai Sign Language: TSL) จำนวนมาก

งานวิจัยนี้นำเสนอระบบ TSL-51 สำหรับการจดจำท่าทางมือในภาษามือไทยจำนวน 51 ท่าทาง

\section{วัตถุประสงค์}

\begin{enumerate}
    \item พัฒนาระบบจดจำภาษามือไทย TSL-51
    \item สร้างส่วนติดต่อผู้ใช้แบบ TUI
    \item พัฒนาระบบแปลวิดีโอเป็นข้อความแบบเรียลไทม์
    \item ทดสอบและเปรียบเทียบประสิทธิภาพของโมเดล
\end{enumerate}

\section{ขอบเขต}

\begin{itemize}
    \item จำนวนท่าทาง: 51 ท่าทาง
    \item ชุดข้อมูล: TSL-51 Dataset
    \item โมเดล: GRU และ MLP
    \item Features: 162 มิติ
\end{itemize}

\chapter{ทฤษฎีและงานวิจัยที่เกี่ยวข้อง}

\section{Deep Learning}

Deep Learning เป็นสาขาหนึ่งของ Machine Learning ที่ใช้ Neural Network หลายชั้น

\section{GRU}

Gated Recurrent Unit (GRU) เป็น variant ของ RNN ที่มีความสามารถในการจดจำลำดับเวลา

\section{MediaPipe}

MediaPipe เป็น Framework ของ Google สำหรับการตรวจจับจุด Landmark บนมือ

\chapter{การออกแบบระบบ}

\section{สถาปัตยกรรม}

ระบบ TSL-51 ประกอบด้วย:
\begin{itemize}
    \item ส่วนการเตรียมข้อมูล
    \item ส่วนการฝึกโมเดล
    \item ส่วนการทำ Inference
    \item ส่วนติดต่อผู้ใช้ TUI
\end{itemize}

\section{โมเดล GRU}

\begin{lstlisting}[language=Python]
class GRUModel(nn.Module):
    def __init__(self, input_dim=162, hidden_dim=256, 
                 num_classes=51, num_layers=3, dropout=0.3):
        super().__init__()
        self.gru = nn.GRU(
            input_size=input_dim,
            hidden_size=hidden_dim,
            num_layers=num_layers,
            batch_first=True,
            bidirectional=True,
        )
        self.fc = nn.Linear(hidden_dim * 2, num_classes)
\end{lstlisting}

\chapter{การทดลองและผลลัพธ์}

\section{ผลการเปรียบเทียบโมเดล}

\begin{longtable}{|p{3cm}|p{3cm}|p{3cm}|}
\hline
โมเดล & Accuracy & F1-Score \\
\hline
GRU (3 layers) & 92.5\% & 0.91 \\
\hline
MLP (3 layers) & 87.3\% & 0.86 \\
\hline
\end{longtable}

\section{การวิเคราะห์}

GRU ให้ผลลัพธ์ที่ดีกว่า MLP เนื่องจากสามารถจับข้อมูลลำดับเวลาได้ดีกว่า

\chapter{บทสรุป}

\section{สรุปผลการวิจัย}

งานวิจัยนี้ได้พัฒนาระบบ TSL-51 สำเร็จ โดยมีผลการวิจัยดังนี้:
\begin{itemize}
    \item ความแม่นยำ 92.5\% ด้วยโมเดล GRU
    \item พัฒนา TUI สำหรับการใช้งาน
    \item รองรับ Real-time Camera Translation
\end{itemize}

\section{ข้อเสนอแนะ}

\begin{enumerate}
    \item เพิ่มจำนวนท่าทางให้ครอบคลุมมากขึ้น
    \item พัฒนา Sentence-level Translation
    \item ปรับปรุง Real-time Performance
\end{enumerate}

\begin{thebibliography}{99}

bibitem{mediapipe}
Google. (2024). MediaPipe. \url{https://mediapipe.dev}

bibitem{tsldataset}
Namonpas. (2024). Thai Sign Language TSL-51 Dataset. HuggingFace.

bibitem{gru}
Cho, K., et al. (2014). Learning Phrase Representations using RNN Encoder-Decoder. EMNLP 2014.

bibitem{pytorch}
Paszke, A., et al. (2019). PyTorch: An Imperative Style, High-Performance Deep Learning Library. NeurIPS 2019.

\end{thebibliography}

\end{document}